\scanpage{400}\mspage{124}
ont assez de cofibrations et assez de fibrations,
donc qu'ils satisfont $\mathrm{Loc}(6\mathrm{a})$ et $\mathrm{Loc}(7\mathrm{F})$.
(Pour ce qui est de $\mathrm{Loc}(6\mathrm{a})$, il semble qu'on
puisse le prouver directement \inserted{très simplement}, cf \S{} suivant.)
Mais ici encore, il faudrait regarder quelques
exemples (équivalences homologiques ou cohomologiques,
$\mathbb{W}_{\unsure{n|m}}$ i.e. les $n$-équiv. d'homotopie.
et regarder déjà si $\mathrm{Loc}(7\mathrm{F})$ est satisfait.

Il ne \deleted{\unsure{sem}} semble pas que le th. de
Thomason ait été \inserted{jusqu'à présent} ou doive être très utile,
\deleted{\unsure{plutôt|théori}} techniquement parlant. Psychologiquement,
je crois qu'il a été important,
pour donner confiance que $\mathrm{Cat}$ est vraiment
une excellente catégorie de modèles
pour $\mathrm{Hot}$. J'ai été vraiment enchanté et
stimulé d'apprendre ce théorème. Mais
la notion d'application de Dwyer (que j'ai
apprise \inserted{en 1983} dans l'article de Thomason) a
été, techniquement, beaucoup plus \unsure{importante|important}
pour moi : c'était exactement ce qui
manquait à ma panoplie de notions et d'\unsure{intuitions}
``homotopiques'' dans $\mathrm{Cat}$ ....
% archive sheet number: 399
